\chapter{Unit 3: Higher Order Differential Equations and Variation of Parameters}

\section{Section A: Multiple Choice Practice Questions (60 Curated Problems)}

\noindent\textbf{Q1. For a non-homogeneous linear differential equation represented as f ( D ) y = X , what is the expression for the Particular Integral (P.I.)?} \hfill [Tag: solution of non-homogeneous linear differential equations with constant coefficients using operator method | Easy]
\begin{enumerate}[label=(\alph*),topsep=2pt,itemsep=1pt,partopsep=0pt,parsep=0pt]
    \item {{f ( D ) +}}
    \item {{f '}}
    \item {{f ( D )}}
    \item {{f ( D ) X}}
\end{enumerate}
\vspace{2mm}

\noindent\textbf{Q2. When finding the Particular Integral for f ( D ) y = e a x , what is the general rule if f ( a )  = 0 ?} \hfill [Tag: solution of non-homogeneous linear differential equations with constant coefficients using operator method | Easy]
\begin{enumerate}[label=(\alph*),topsep=2pt,itemsep=1pt,partopsep=0pt,parsep=0pt]
    \item {{Replace D 2}}
    \item {{Replace D}}
    \item {{Replace D}}
    \item {{Replace D}}
\end{enumerate}
\vspace{2mm}

\noindent\textbf{Q3. What is the general solution of a non-homogeneous linear differential equation?} \hfill [Tag: solution of non-homogeneous linear differential equations with constant coefficients using operator method | Easy]
\begin{enumerate}[label=(\alph*),topsep=2pt,itemsep=1pt,partopsep=0pt,parsep=0pt]
    \item {{Particular Integral only}}
    \item {{Complementary Function + Particular Integral}}
    \item {{Complementary Function $\times$}}
    \item {{Complementary Function only}}
\end{enumerate}
\vspace{2mm}

\noindent\textbf{Q4. When finding the P.I. for f ( D ) y = sin ( a x ) or cos ( a x ) , what substitution is made if f ( - a 2 )  = 0 ?} \hfill [Tag: solution of non-homogeneous linear differential equations with constant coefficients using operator method | Easy]
\begin{enumerate}[label=(\alph*),topsep=2pt,itemsep=1pt,partopsep=0pt,parsep=0pt]
    \item {{Replace D 2}}
    \item {{Replace D}}
    \item {{Replace D}}
    \item {{Replace D 2}}
\end{enumerate}
\vspace{2mm}

\noindent\textbf{Q5. If f ( a ) = 0 when calculating the P.I. for f ( D ) y = e a x , what is the next step?} \hfill [Tag: solution of non-homogeneous linear differential equations with constant coefficients using operator method | Easy]
\begin{enumerate}[label=(\alph*),topsep=2pt,itemsep=1pt,partopsep=0pt,parsep=0pt]
    \item {{Divide by x}}
    \item {{The P.I. is zero}}
    \item {{Multiply by x}}
    \item {{The equation has no solution}}
\end{enumerate}
\vspace{2mm}

\noindent\textbf{Q6. What is the primary purpose of the method of variation of parameters?} \hfill [Tag: method of variation of parameters | Easy]
\begin{enumerate}[label=(\alph*),topsep=2pt,itemsep=1pt,partopsep=0pt,parsep=0pt]
    \item {{To find the Particular Integral}}
    \item {{To convert variable coefficients into constant coefficients}}
    \item {{To find the Complementary Function}}
    \item {{To find the roots of the auxiliary equation}}
\end{enumerate}
\vspace{2mm}

\noindent\textbf{Q7. In the method of variation of parameters for a second-order equation, what is the Wronskian W of two linearly independent solutions y 1 and y 2 ?} \hfill [Tag: method of variation of parameters | Easy]
\begin{enumerate}[label=(\alph*),topsep=2pt,itemsep=1pt,partopsep=0pt,parsep=0pt]
    \item {{y 1}}
    \item {{y 1}}
    \item {{y 1}}
    \item {{y 1}}
\end{enumerate}
\vspace{2mm}

\noindent\textbf{Q8. For two solutions y 1 and y 2 to be linearly independent, what must be true about their Wronskian W ?} \hfill [Tag: method of variation of parameters | Easy]
\begin{enumerate}[label=(\alph*),topsep=2pt,itemsep=1pt,partopsep=0pt,parsep=0pt]
    \item {{W }}
    \item {{W =}}
    \item {{W =}}
    \item {{W =}}
\end{enumerate}
\vspace{2mm}

\noindent\textbf{Q9. In the formula for the Particular Integral y p = u y 1 + v y 2 , what is the standard formula for u ?} \hfill [Tag: method of variation of parameters | Easy]
\begin{enumerate}[label=(\alph*),topsep=2pt,itemsep=1pt,partopsep=0pt,parsep=0pt]
    \item {{$\int$ W}}
    \item {{$\int$ W}}
    \item {{- $\int$ W}}
    \item {{- $\int$ W}}
\end{enumerate}
\vspace{2mm}

\noindent\textbf{Q10. Which of the following makes the method of variation of parameters more powerful than the undetermined coefficients method?} \hfill [Tag: method of variation of parameters | Easy]
\begin{enumerate}[label=(\alph*),topsep=2pt,itemsep=1pt,partopsep=0pt,parsep=0pt]
    \item {{It does not require integration}}
    \item {{It applies only to constant coefficients}}
    \item {{It applies even when the right-hand side is not a simple exponential, polynomial, or sine/cosine function}}
    \item {{It requires fewer calculations}}
\end{enumerate}
\vspace{2mm}

\noindent\textbf{Q11. The method of undetermined coefficients is best suited for right-hand side functions that are:} \hfill [Tag: method of undetermined coefficient | Easy]
\begin{enumerate}[label=(\alph*),topsep=2pt,itemsep=1pt,partopsep=0pt,parsep=0pt]
    \item {{Inverse trigonometric functions}}
    \item {{Logarithmic functions}}
    \item {{Polynomials, exponentials, sines, and cosines}}
    \item {{Fractional functions like 1/ x}}
\end{enumerate}
\vspace{2mm}

\noindent\textbf{Q12. Can the method of undetermined coefficients be used to find the particular integral if the right-hand side is tan ( x ) ?} \hfill [Tag: method of undetermined coefficient | Easy]
\begin{enumerate}[label=(\alph*),topsep=2pt,itemsep=1pt,partopsep=0pt,parsep=0pt]
    \item {{Yes, but only if the degree is 1}}
    \item {{Yes, by converting it to sin ( x )}}
    \item {{No, it is not applicable}}
    \item {{Yes, always}}
\end{enumerate}
\vspace{2mm}

\noindent\textbf{Q13. If the non-homogeneous term is X = e 2 x and 2 is NOT a root of the auxiliary equation, what is the appropriate trial solution?} \hfill [Tag: method of undetermined coefficient | Easy]
\begin{enumerate}[label=(\alph*),topsep=2pt,itemsep=1pt,partopsep=0pt,parsep=0pt]
    \item {{A e 2 x}}
    \item {{A x 2}}
    \item {{A e - 2 x}}
    \item {{A x e 2 x}}
\end{enumerate}
\vspace{2mm}

\noindent\textbf{Q14. If the non-homogeneous term X is a polynomial of degree n , what should be the form of the trial solution?} \hfill [Tag: method of undetermined coefficient | Easy]
\begin{enumerate}[label=(\alph*),topsep=2pt,itemsep=1pt,partopsep=0pt,parsep=0pt]
    \item {{A constant}}
    \item {{A trigonometric function}}
    \item {{An exponential function}}
    \item {{A polynomial of degree n}}
\end{enumerate}
\vspace{2mm}

\noindent\textbf{Q15. After substituting the trial solution into the differential equation, how are the undetermined coefficients found?} \hfill [Tag: method of undetermined coefficient | Easy]
\begin{enumerate}[label=(\alph*),topsep=2pt,itemsep=1pt,partopsep=0pt,parsep=0pt]
    \item {{By equating the coefficients of similar terms on both sides}}
    \item {{By integrating both sides}}
    \item {{By setting the equation to zero}}
    \item {{By taking the Laplace transform}}
\end{enumerate}
\vspace{2mm}

\noindent\textbf{Q16. What is the standard form of a second-order Euler-Cauchy differential equation?} \hfill [Tag: solution of Euler-Cauchy equation | Easy]
\begin{enumerate}[label=(\alph*),topsep=2pt,itemsep=1pt,partopsep=0pt,parsep=0pt]
    \item {{x 2}}
    \item {{x y ''}}
    \item {{y ''}}
    \item {{y ''}}
\end{enumerate}
\vspace{2mm}

\noindent\textbf{Q17. Which standard substitution is used to transform an Euler-Cauchy equation into a linear differential equation with constant coefficients?} \hfill [Tag: solution of Euler-Cauchy equation | Easy]
\begin{enumerate}[label=(\alph*),topsep=2pt,itemsep=1pt,partopsep=0pt,parsep=0pt]
    \item {{x =}}
    \item {{x =}}
    \item {{x =}}
    \item {{x =}}
\end{enumerate}
\vspace{2mm}

\noindent\textbf{Q18. Using the substitution x = e t and letting $\theta$ = d t d , what does the operator x d x d reduce to?} \hfill [Tag: solution of Euler-Cauchy equation | Easy]
\begin{enumerate}[label=(\alph*),topsep=2pt,itemsep=1pt,partopsep=0pt,parsep=0pt]
    \item {{$\theta$ 2}}
    \item {{$\theta$ +}}
    \item {{$\theta$ -}}
    \item {{$\theta$}}
\end{enumerate}
\vspace{2mm}

\noindent\textbf{Q19. Under the standard transformation x = e t with $\theta$ = d t d , what is the equivalent operator for x 2 d x 2 d 2 ?} \hfill [Tag: solution of Euler-Cauchy equation | Easy]
\begin{enumerate}[label=(\alph*),topsep=2pt,itemsep=1pt,partopsep=0pt,parsep=0pt]
    \item {{2 $\theta$}}
    \item {{$\theta$ 2}}
    \item {{$\theta$ ( $\theta$ -}}
    \item {{$\theta$ ( $\theta$ +}}
\end{enumerate}
\vspace{2mm}

\noindent\textbf{Q20. An Euler-Cauchy differential equation is an example of a linear differential equation with:} \hfill [Tag: solution of Euler-Cauchy equation | Easy]
\begin{enumerate}[label=(\alph*),topsep=2pt,itemsep=1pt,partopsep=0pt,parsep=0pt]
    \item {{Only trigonometric coefficients}}
    \item {{Constant coefficients}}
    \item {{No coefficients}}
    \item {{Variable coefficients}}
\end{enumerate}
\vspace{2mm}

\noindent\textbf{Q21. Find the Particular Integral (P.I.) of the differential equation ( D 2 - 4 ) y = e 2 x .} \hfill [Tag: solution of non-homogeneous linear differential equations with constant coefficients using operator method | Medium]
\begin{enumerate}[label=(\alph*),topsep=2pt,itemsep=1pt,partopsep=0pt,parsep=0pt]
    \item {{x e 2 x}}
    \item {{4}}
    \item {{4}}
    \item {{2}}
\end{enumerate}
\vspace{2mm}

\noindent\textbf{Q22. What is the Particular Integral for ( D 2 + 1 ) y = sin x ?} \hfill [Tag: solution of non-homogeneous linear differential equations with constant coefficients using operator method | Medium]
\begin{enumerate}[label=(\alph*),topsep=2pt,itemsep=1pt,partopsep=0pt,parsep=0pt]
    \item {{- 2}}
    \item {{2}}
    \item {{2}}
    \item {{- 2}}
\end{enumerate}
\vspace{2mm}

\noindent\textbf{Q23. Find the Particular Integral of ( D 2 - 5 D + 6 ) y = x 2 .} \hfill [Tag: solution of non-homogeneous linear differential equations with constant coefficients using operator method | Medium]
\begin{enumerate}[label=(\alph*),topsep=2pt,itemsep=1pt,partopsep=0pt,parsep=0pt]
    \item {{6}}
    \item {{6}}
    \item {{3}}
    \item {{6}}
\end{enumerate}
\vspace{2mm}

\noindent\textbf{Q24. The P.I. of ( D 2 - 2 D + 1 ) y = x e x is:} \hfill [Tag: solution of non-homogeneous linear differential equations with constant coefficients using operator method | Medium]
\begin{enumerate}[label=(\alph*),topsep=2pt,itemsep=1pt,partopsep=0pt,parsep=0pt]
    \item {{6}}
    \item {{2}}
    \item {{12}}
    \item {{3}}
\end{enumerate}
\vspace{2mm}

\noindent\textbf{Q25. In the method of variation of parameters for y '' + p ( x ) y ' + q ( x ) y = R ( x ) , if y 1 and y 2 are independent solutions of the homogeneous part, the Wronskian W is given by:} \hfill [Tag: method of variation of parameters | Medium]
\begin{enumerate}[label=(\alph*),topsep=2pt,itemsep=1pt,partopsep=0pt,parsep=0pt]
    \item {{y 1}}
    \item {{y 1}}
    \item {{y 1}}
    \item {{y 1}}
\end{enumerate}
\vspace{2mm}

\noindent\textbf{Q26. For the equation y '' + y = sec x , the Wronskian of the fundamental solutions y 1 = cos x and y 2 = sin x is:} \hfill [Tag: method of variation of parameters | Medium]
\begin{enumerate}[label=(\alph*),topsep=2pt,itemsep=1pt,partopsep=0pt,parsep=0pt]
    \item {{1}}
    \item {{sin 2 x}}
    \item {{- 1}}
    \item {{cos 2 x}}
\end{enumerate}
\vspace{2mm}

\noindent\textbf{Q27. Using variation of parameters for y '' + y = sec x , the Particular Integral is u ( x ) y 1 + v ( x ) y 2 . What is u ( x ) if y 1 = cos x and y 2 = sin x ?} \hfill [Tag: method of variation of parameters | Medium]
\begin{enumerate}[label=(\alph*),topsep=2pt,itemsep=1pt,partopsep=0pt,parsep=0pt]
    \item {{ln ( sin x )}}
    \item {{ln ( cos x )}}
    \item {{sec x}}
    \item {{x}}
\end{enumerate}
\vspace{2mm}

\noindent\textbf{Q28. Using variation of parameters for y '' + 4 y = tan 2 x , what is the Wronskian W of y 1 = cos 2 x and y 2 = sin 2 x ?} \hfill [Tag: method of variation of parameters | Medium]
\begin{enumerate}[label=(\alph*),topsep=2pt,itemsep=1pt,partopsep=0pt,parsep=0pt]
    \item {{- 2}}
    \item {{4}}
    \item {{2}}
    \item {{1}}
\end{enumerate}
\vspace{2mm}

\noindent\textbf{Q29. In the method of undetermined coefficients for solving y '' - 3 y ' + 2 y = e 3 x , the assumed form of the Particular Integral is:} \hfill [Tag: method of undetermined coefficient | Medium]
\begin{enumerate}[label=(\alph*),topsep=2pt,itemsep=1pt,partopsep=0pt,parsep=0pt]
    \item {{A e 3 x}}
    \item {{A e 2 x}}
    \item {{A x 2}}
    \item {{A x e 3 x}}
\end{enumerate}
\vspace{2mm}

\noindent\textbf{Q30. To find the P.I. for y '' - 4 y = e 2 x using the method of undetermined coefficients, the appropriate trial function is:} \hfill [Tag: method of undetermined coefficient | Medium]
\begin{enumerate}[label=(\alph*),topsep=2pt,itemsep=1pt,partopsep=0pt,parsep=0pt]
    \item {{A e 2 x}}
    \item {{A x 2}}
    \item {{A x e 2 x}}
    \item {{A e - 2 x}}
\end{enumerate}
\vspace{2mm}

\noindent\textbf{Q31. What is the correct trial solution for the Particular Integral of y '' + y = x sin x using the method of undetermined coefficients?} \hfill [Tag: method of undetermined coefficient | Medium]
\begin{enumerate}[label=(\alph*),topsep=2pt,itemsep=1pt,partopsep=0pt,parsep=0pt]
    \item {{x 2}}
    \item {{( A x +}}
    \item {{x (( A x +}}
    \item {{A x sin x +}}
\end{enumerate}
\vspace{2mm}

\noindent\textbf{Q32. Determine the undetermined coefficient A if the P.I. for y '' + 3 y ' + 2 y = 4 x is assumed to be A x + B .} \hfill [Tag: method of undetermined coefficient | Medium]
\begin{enumerate}[label=(\alph*),topsep=2pt,itemsep=1pt,partopsep=0pt,parsep=0pt]
    \item {{1}}
    \item {{4}}
    \item {{- 3}}
    \item {{2}}
\end{enumerate}
\vspace{2mm}

\noindent\textbf{Q33. To solve the Euler-Cauchy equation x 2 y '' - 2 x y ' + 2 y = 0 , which substitution is generally used to convert it to a linear equation with constant coefficients?} \hfill [Tag: solution of Euler-Cauchy equation | Medium]
\begin{enumerate}[label=(\alph*),topsep=2pt,itemsep=1pt,partopsep=0pt,parsep=0pt]
    \item {{y =}}
    \item {{x =}}
    \item {{x =}}
    \item {{y =}}
\end{enumerate}
\vspace{2mm}

\noindent\textbf{Q34. Under the substitution x = e z (where D = d z d ), what does the term x 2 d x 2 d 2 y become?} \hfill [Tag: solution of Euler-Cauchy equation | Medium]
\begin{enumerate}[label=(\alph*),topsep=2pt,itemsep=1pt,partopsep=0pt,parsep=0pt]
    \item {{D 2}}
    \item {{( D 2}}
    \item {{D ( D -}}
    \item {{D ( D +}}
\end{enumerate}
\vspace{2mm}

\noindent\textbf{Q35. Solve the Euler-Cauchy equation x 2 y '' - 3 x y ' + 4 y = 0 .} \hfill [Tag: solution of Euler-Cauchy equation | Medium]
\begin{enumerate}[label=(\alph*),topsep=2pt,itemsep=1pt,partopsep=0pt,parsep=0pt]
    \item {{y =}}
    \item {{y =}}
    \item {{y =}}
    \item {{y =}}
\end{enumerate}
\vspace{2mm}

\noindent\textbf{Q36. Find the general solution to x 2 y '' + 5 x y ' + 4 y = 0 .} \hfill [Tag: solution of Euler-Cauchy equation | Medium]
\begin{enumerate}[label=(\alph*),topsep=2pt,itemsep=1pt,partopsep=0pt,parsep=0pt]
    \item {{y =}}
    \item {{y =}}
    \item {{y =}}
    \item {{y =}}
\end{enumerate}
\vspace{2mm}

\noindent\textbf{Q37. For the non-homogeneous Euler-Cauchy equation x 2 y '' - x y ' = x 3 , what is the differential equation in terms of z after substitution x = e z ?} \hfill [Tag: solution of Euler-Cauchy equation | Medium]
\begin{enumerate}[label=(\alph*),topsep=2pt,itemsep=1pt,partopsep=0pt,parsep=0pt]
    \item {{( D 2}}
    \item {{( D 2}}
    \item {{( D 2}}
    \item {{( D 2}}
\end{enumerate}
\vspace{2mm}

\noindent\textbf{Q38. Find the Particular Integral of ( D 2 - 1 ) y = cosh x .} \hfill [Tag: solution of non-homogeneous linear differential equations with constant coefficients using operator method | Medium]
\begin{enumerate}[label=(\alph*),topsep=2pt,itemsep=1pt,partopsep=0pt,parsep=0pt]
    \item {{2}}
    \item {{2}}
    \item {{2}}
    \item {{x cosh x}}
\end{enumerate}
\vspace{2mm}

\noindent\textbf{Q39. In variation of parameters for y '' + y = csc x , find the function v ( x ) given P . I . = u ( x ) cos x + v ( x ) sin x .} \hfill [Tag: method of variation of parameters | Medium]
\begin{enumerate}[label=(\alph*),topsep=2pt,itemsep=1pt,partopsep=0pt,parsep=0pt]
    \item {{ln ( cos x )}}
    \item {{- cos x}}
    \item {{ln ( sin x )}}
    \item {{x}}
\end{enumerate}
\vspace{2mm}

\noindent\textbf{Q40. If R ( x ) = x 2 e x , and the roots of the characteristic equation are m = 1 , 2 , what is the assumed form for the Particular Integral in the method of undetermined coefficients?} \hfill [Tag: method of undetermined coefficient | Medium]
\begin{enumerate}[label=(\alph*),topsep=2pt,itemsep=1pt,partopsep=0pt,parsep=0pt]
    \item {{x 2}}
    \item {{x ( A x 2}}
    \item {{( A x 2}}
    \item {{( A x +}}
\end{enumerate}
\vspace{2mm}

\noindent\textbf{Q41. Find the particular integral of ( D 2 + 4 ) y = x sin 2 x .} \hfill [Tag: solution of non-homogeneous linear differential equations with constant coefficients using operator method | Hard]
\begin{enumerate}[label=(\alph*),topsep=2pt,itemsep=1pt,partopsep=0pt,parsep=0pt]
    \item {{- 8}}
    \item {{8}}
    \item {{- 4}}
    \item {{4}}
\end{enumerate}
\vspace{2mm}

\noindent\textbf{Q42. In the method of variation of parameters for y '' + y = sec 3 x , what are the derivatives of the parameters v 1 ' and v 2 ' ?} \hfill [Tag: method of variation of parameters | Hard]
\begin{enumerate}[label=(\alph*),topsep=2pt,itemsep=1pt,partopsep=0pt,parsep=0pt]
    \item {{v 1}}
    \item {{v 1}}
    \item {{v 1}}
    \item {{v 1}}
\end{enumerate}
\vspace{2mm}

\noindent\textbf{Q43. For the differential equation y '' - 4 y ' + 4 y = x 2 e 2 x , what is the correct assumed form for the particular solution using the method of undetermined coefficients?} \hfill [Tag: method of undetermined coefficient | Hard]
\begin{enumerate}[label=(\alph*),topsep=2pt,itemsep=1pt,partopsep=0pt,parsep=0pt]
    \item {{y p}}
    \item {{y p}}
    \item {{y p}}
    \item {{y p}}
\end{enumerate}
\vspace{2mm}

\noindent\textbf{Q44. Find the general solution of the Euler-Cauchy equation x 3 y ''' + x 2 y '' - 2 x y ' + 2 y = 0 .} \hfill [Tag: solution of Euler-Cauchy equation | Hard]
\begin{enumerate}[label=(\alph*),topsep=2pt,itemsep=1pt,partopsep=0pt,parsep=0pt]
    \item {{y =}}
    \item {{y =}}
    \item {{y =}}
    \item {{y =}}
\end{enumerate}
\vspace{2mm}

\noindent\textbf{Q45. What is the particular integral of ( D 2 + 1 ) 2 y = x 2 cos x ?} \hfill [Tag: solution of non-homogeneous linear differential equations with constant coefficients using operator method | Hard]
\begin{enumerate}[label=(\alph*),topsep=2pt,itemsep=1pt,partopsep=0pt,parsep=0pt]
    \item {{24}}
    \item {{48}}
    \item {{- 24}}
    \item {{- 48}}
\end{enumerate}
\vspace{2mm}

\noindent\textbf{Q46. Using the method of variation of parameters for the equation x 2 y '' - 2 x y ' + 2 y = x 4 e x , what is the Wronskian W ( y 1 , y 2 ) of the complementary solutions if evaluated correctly as functions of x ?} \hfill [Tag: method of variation of parameters | Hard]
\begin{enumerate}[label=(\alph*),topsep=2pt,itemsep=1pt,partopsep=0pt,parsep=0pt]
    \item {{e x}}
    \item {{2 x}}
    \item {{x 4}}
    \item {{x 2}}
\end{enumerate}
\vspace{2mm}

\noindent\textbf{Q47. For the equation y ( 4 ) - y = cos x + x e - x , what is the minimum number of undetermined coefficients needed to find the particular solution?} \hfill [Tag: method of undetermined coefficient | Hard]
\begin{enumerate}[label=(\alph*),topsep=2pt,itemsep=1pt,partopsep=0pt,parsep=0pt]
    \item {{5}}
    \item {{6}}
    \item {{4}}
    \item {{7}}
\end{enumerate}
\vspace{2mm}

\noindent\textbf{Q48. Consider the non-homogeneous Euler-Cauchy equation x 2 y '' + x y ' - y = x ln x . Under the substitution x = e t , what is the corresponding differential equation in t ?} \hfill [Tag: solution of Euler-Cauchy equation | Hard]
\begin{enumerate}[label=(\alph*),topsep=2pt,itemsep=1pt,partopsep=0pt,parsep=0pt]
    \item {{d t 2}}
    \item {{d t 2}}
    \item {{d t 2}}
    \item {{d t 2}}
\end{enumerate}
\vspace{2mm}

\noindent\textbf{Q49. Find the particular integral of ( D 3 - D 2 - 6 D ) y = 1 + x 2 .} \hfill [Tag: solution of non-homogeneous linear differential equations with constant coefficients using operator method | Hard]
\begin{enumerate}[label=(\alph*),topsep=2pt,itemsep=1pt,partopsep=0pt,parsep=0pt]
    \item {{6}}
    \item {{- 6}}
    \item {{- 6}}
    \item {{- 6}}
\end{enumerate}
\vspace{2mm}

\noindent\textbf{Q50. For the differential equation y '' + 4 y = tan 2 x , the particular solution obtained using variation of parameters is:} \hfill [Tag: method of variation of parameters | Hard]
\begin{enumerate}[label=(\alph*),topsep=2pt,itemsep=1pt,partopsep=0pt,parsep=0pt]
    \item {{- 4}}
    \item {{2}}
    \item {{4}}
    \item {{- 2}}
\end{enumerate}
\vspace{2mm}

\noindent\textbf{Q51. When solving y ''' - y '' = 3 e x + x by the method of undetermined coefficients, what is the appropriate form of the particular solution?} \hfill [Tag: method of undetermined coefficient | Hard]
\begin{enumerate}[label=(\alph*),topsep=2pt,itemsep=1pt,partopsep=0pt,parsep=0pt]
    \item {{y p}}
    \item {{y p}}
    \item {{y p}}
    \item {{y p}}
\end{enumerate}
\vspace{2mm}

\noindent\textbf{Q52. Find the general solution of ( 2 x + 3 ) 2 y '' - 2 ( 2 x + 3 ) y ' - 12 y = 6 x .} \hfill [Tag: solution of Euler-Cauchy equation | Hard]
\begin{enumerate}[label=(\alph*),topsep=2pt,itemsep=1pt,partopsep=0pt,parsep=0pt]
    \item {{y =}}
    \item {{y =}}
    \item {{y =}}
    \item {{y =}}
\end{enumerate}
\vspace{2mm}

\noindent\textbf{Q53. Evaluate the particular integral for ( D 2 + 2 D + 2 ) y = e - x sec x .} \hfill [Tag: solution of non-homogeneous linear differential equations with constant coefficients using operator method | Hard]
\begin{enumerate}[label=(\alph*),topsep=2pt,itemsep=1pt,partopsep=0pt,parsep=0pt]
    \item {{e - x}}
    \item {{x e - x}}
    \item {{e - x}}
    \item {{e - x}}
\end{enumerate}
\vspace{2mm}

\noindent\textbf{Q54. Apply variation of parameters to y '' - 2 y ' + y = x 2 + 1 e x . The particular solution is:} \hfill [Tag: method of variation of parameters | Hard]
\begin{enumerate}[label=(\alph*),topsep=2pt,itemsep=1pt,partopsep=0pt,parsep=0pt]
    \item {{y p}}
    \item {{y p}}
    \item {{y p}}
    \item {{y p}}
\end{enumerate}
\vspace{2mm}

\noindent\textbf{Q55. For the system corresponding to the differential equation y '' - 2 y ' + 5 y = 8 e x cos 2 x , what is the correct form for the particular solution?} \hfill [Tag: method of undetermined coefficient | Hard]
\begin{enumerate}[label=(\alph*),topsep=2pt,itemsep=1pt,partopsep=0pt,parsep=0pt]
    \item {{y p}}
    \item {{y p}}
    \item {{y p}}
    \item {{y p}}
\end{enumerate}
\vspace{2mm}

\noindent\textbf{Q56. Solve the initial value problem for the Euler-Cauchy equation x 2 y '' - 3 x y ' + 4 y = 0 with y ( 1 ) = 2 and y ' ( 1 ) = 5 .} \hfill [Tag: solution of Euler-Cauchy equation | Hard]
\begin{enumerate}[label=(\alph*),topsep=2pt,itemsep=1pt,partopsep=0pt,parsep=0pt]
    \item {{y ( x ) =}}
    \item {{y ( x ) =}}
    \item {{y ( x ) =}}
    \item {{y ( x ) =}}
\end{enumerate}
\vspace{2mm}

\noindent\textbf{Q57. Find the particular integral of ( D 2 + a 2 ) y = sec a x .} \hfill [Tag: solution of non-homogeneous linear differential equations with constant coefficients using operator method | Hard]
\begin{enumerate}[label=(\alph*),topsep=2pt,itemsep=1pt,partopsep=0pt,parsep=0pt]
    \item {{a}}
    \item {{a}}
    \item {{2 a}}
    \item {{a}}
\end{enumerate}
\vspace{2mm}

\noindent\textbf{Q58. Which of the following represents the Wronskian W of the fundamental solutions of the differential equation ( 1 - x 2 ) y '' - 2 x y ' + 2 y = 0 , given that one solution is y 1 = x ?} \hfill [Tag: method of variation of parameters | Hard]
\begin{enumerate}[label=(\alph*),topsep=2pt,itemsep=1pt,partopsep=0pt,parsep=0pt]
    \item {{C e - x 2}}
    \item {{C ( 1 -}}
    \item {{1 - x 2}}
    \item {{( 1 - x 2}}
\end{enumerate}
\vspace{2mm}

\noindent\textbf{Q59. If the annihilator method is used to find the particular solution form for y '' - y = x e x sin x , what is the minimum order of the annihilator operator?} \hfill [Tag: method of undetermined coefficient | Hard]
\begin{enumerate}[label=(\alph*),topsep=2pt,itemsep=1pt,partopsep=0pt,parsep=0pt]
    \item {{2}}
    \item {{4}}
    \item {{3}}
    \item {{6}}
\end{enumerate}
\vspace{2mm}

\noindent\textbf{Q60. Solve the Euler-Cauchy equation x 2 y '' + 5 x y ' + 4 y = x - 2 ln x .} \hfill [Tag: solution of Euler-Cauchy equation | Hard]
\begin{enumerate}[label=(\alph*),topsep=2pt,itemsep=1pt,partopsep=0pt,parsep=0pt]
    \item {{y =}}
    \item {{y =}}
    \item {{y =}}
    \item {{y =}}
\end{enumerate}
\vspace{2mm}


\section{Section B: Comprehensive Subjective Problem with Detailed Analytical Rubric}

\begin{subjectiveq}{3}{1}
Solve by the method of Variation of Parameters: $\frac{d^2 y}{dx^2} + y = \sec x$.
\end{subjectiveq}

\begin{solbox}
Complementary function: $m^2 + 1 = 0 \implies y_c = c_1 \cos x + c_2 \sin x$. Here $y_1 = \cos x$, $y_2 = \sin x$. The Wronskian is $W = y_1 y_2' - y_2 y_1' = \cos^2 x + \sin^2 x = 1$. The particular integral is $y_p = u(x)y_1 + v(x)y_2$, where $u(x) = -\int \frac{\sin x \sec x}{1} dx = -\int \tan x dx = \ln|\cos x|$, and $v(x) = \int \frac{\cos x \sec x}{1} dx = \int 1 dx = x$. General solution: $y = c_1 \cos x + c_2 \sin x + \cos x \ln|\cos x| + x \sin x$.
\end{solbox}

\begin{examinertip}
Step 1: Identify the foundational definition and state governing theorems.\\
Step 2: Clearly write each algebraic transformation and boundary substitution.\\
Step 3: State the final numerical/algebraic invariant clearly in a concluding box.
\end{examinertip}
